\documentclass[a4paper]{article}

\usepackage[margin=2cm]{geometry}
\usepackage{amssymb}
\usepackage{pdfpages}
\usepackage{graphbox}
\usepackage{tikz}
\begin{document}

\begin{center}
  {\Large\bf Percentiles of the chi-square distribution}
\end{center}

\rule{0mm}{5mm}
\pagestyle{empty}
\thispagestyle{empty}
\begin{table}[ht]
\centering
\begin{tabular}{r|llllllllll}
  \begin{minipage}{20mm}degrees\\of freedom\end{minipage}
  & {\Large $\chi^2_{0.9}$}
  & {\Large $\chi^2_{0.95}$}
  & {\Large $\chi^2_{0.975}$}
  & {\Large $\chi^2_{0.99}$}
  & {\Large $\chi^2_{0.995}$}
  \rule[-4mm]{0mm}{10mm}
  \\  \hline
\input{chisq_table_values.txt}
  \hline
\end{tabular}
\end{table}

\begin{centering}
\includegraphics[align=c,width=2.2in]{chisq_diagram.pdf}
\end{centering}

\noindent For example, to calculate the~$95^\mathrm{th}$ percentile of
the chi-square distribution with~30 degrees of freedom, the last line
gives the different quantiles.  Choose the entry with column
heading~$\chi^2_{0.95}$ to obtain~43.8.
\\
\begin{tikzpicture}[remember picture,overlay]
  \node[anchor=south east,
    xshift=-1cm,
    yshift=1cm]
  at (current page.south east){
    {\tt https://github.com/RobinHankin/tables.git}\rule{90mm}{0mm}\includegraphics[width=1in]{AUT-logo-block.pdf}};
\end{tikzpicture}


\end{document}
